% ============================================================
%  CAPITOLO 1 --- INTRODUZIONE
% ============================================================
\chapter{Introduzione}

Il presente elaborato descrive la progettazione e l'implementazione di un sistema
informativo per la gestione di uno \textbf{stabilimento balneare}.

L'obiettivo del progetto è realizzare un gestionale che i proprietari dello stabilimento
possano utilizzare per il monitoraggio delle funzioni essenziali: prenotazioni di ombrelloni
e lettini, abbonamenti stagionali, gestione dei turni del personale e controllo del magazzino
e delle forniture.

\section{Contesto e motivazioni}

La gestione di uno stabilimento balneare coinvolge diverse aree operative che, se
non coordinate, rischiano di generare inefficienze: l'assegnazione manuale degli
ombrelloni può portare a doppie prenotazioni, la mancanza di uno strumento per i
turni rende difficile pianificare il personale, e l'assenza di controllo sul magazzino
causa carenze di materiali.

Il sistema informativo proposto affronta queste criticità fornendo al proprietario
un unico strumento di consultazione e gestione, in modo da rendere più efficienti alcune azioni 
che, come prima descritto, potrebbero portare a disservizi nello stabilimento.

\section{Ambito del sistema}

Il progetto \textbf{non} si occupa delle interazioni tra clienti e stabilimento:
non è un sistema di prenotazione online accessibile dagli utenti finali. Le informazioni
trattate sono esclusivamente rivolte ai \textbf{proprietari}, che accedono al gestionale
tramite credenziali personali.

\section{Funzionalità principali}

Il sistema offre le seguenti funzionalità:

\begin{itemize}
  \item Inserimento e gestione di \textbf{clienti} (abituali o occasionali)
  \item Gestione di \textbf{prenotazioni} e \textbf{abbonamenti} su ombrelloni e lettini
  \item Registrazione e gestione dei \textbf{dipendenti} e dei loro \textbf{turni}
  \item Monitoraggio del \textbf{magazzino} e gestione degli \textbf{ordini} ai fornitori
  \item Consultazione della disponibilità in tempo reale
  \item Controllo dello storico di prenotazioni, abbonamenti, ordini e turni lavorativi
\end{itemize}

\section{Struttura della relazione}

La relazione è organizzata come segue:

\begin{description}
  \item[Capitolo 2 --- Analisi dei requisiti:] intervista, rilevamento ambiguità,
    specifiche in linguaggio naturale ed estrazione dei concetti principali.
  \item[Capitolo 3 --- Progettazione concettuale:] schema scheletro, raffinamenti
    progressivi e schema E/R finale.
  \item[Capitolo 4 --- Progettazione logica:] stima dei volumi, operazioni principali,
    schemi di navigazione, tabelle degli accessi, analisi delle ridondanze, schema
    relazionale finale e query SQL.
  \item[Capitolo 5 --- Progettazione dell'applicazione:] architettura del software,
    tecnologie utilizzate e screenshot dell'interfaccia.
\end{description}
